% Diagrama: Proceso 1 de Apimetro — ETL (Extraction, Treatment and Liberation)
% Dos fuentes (datos abiertos + OSM) → ETL en Python (dos fases) → PostgreSQL/PostGIS
% Basado en: TesisUrbanismo 4-1-1-Intro.tex (flujo) y Anexo A2-ETL.tex (fases)
% Usado en: secciones/04-apimetro.tex
\begin{tikzpicture}[
  fuente/.style = {rectangle, rounded corners=5pt,
                   draw=SotmVerde, line width=0.8pt,
                   fill=SotmCrema, text width=4.8cm,
                   align=flush left, minimum height=1.1cm, font=\scriptsize},
  fase/.style   = {rectangle, rounded corners=4pt,
                   draw=SotmVerde!70, fill=white, text width=4.4cm,
                   align=flush left, minimum height=0.8cm, font=\scriptsize},
  marco/.style  = {rectangle, rounded corners=6pt,
                   draw=SotmVerde, line width=0.8pt, dashed,
                   fill=SotmOlivo!12, inner sep=0.25cm},
  result/.style = {rectangle, rounded corners=5pt,
                   fill=SotmTerracota, draw=SotmTerracota, text=white,
                   text width=3.4cm, align=center, minimum height=1.4cm,
                   font=\scriptsize},
  flecha/.style = {-Stealth, thick, color=SotmVerde}
]
  % ── Fuentes ──
  \node[fuente] (f1) at (0,  1.0) {{\small\bfseries Datos abiertos}\\
    GTFS (ADIP) · SHP (INEGI)\\ parámetros operativos (SEMOVI)};
  \node[fuente] (f2) at (0, -1.0) {{\small\bfseries OpenStreetMap}\\
    Mexibús · Mexicable\\ Trolebús Línea 13};

  % ── ETL: dos fases ──
  \node[fase] (p1) at (6.8,  0.5) {\textbf{1.} GTFS → plantilla con campos
    normalizados (líneas, ramales, estaciones)};
  \node[fase] (p2) at (6.8, -0.8) {\textbf{2.} Cruce de geometrías y carga
    a la base de datos};
  \draw[-Stealth, color=SotmVerde!70] (p1) -- (p2);
  \begin{scope}[on background layer]
    \node[marco, fit=(p1)(p2), inner ysep=0.35cm,
          label={[font=\scriptsize\bfseries, text=SotmVerde, yshift=-0.1cm]above:%
          ETL · \textit{Extraction, Treatment and Liberation}}] (etl) {};
  \end{scope}
  \node[font=\tiny\itshape, text=SotmVerde, anchor=north] at (etl.south)
    {Python · GeoPandas};

  % ── Resultado ──
  \node[result] (db) at (12.6, 0) {{\small\bfseries PostgreSQL + PostGIS}\\[2pt]
    estaciones · líneas\\ polígonos};

  % ── Flujo ──
  \draw[flecha] (f1.east) -- (etl.west |- f1.east);
  \draw[flecha] (f2.east) -- (etl.west |- f2.east);
  \draw[flecha] (etl.east) -- (db.west);
\end{tikzpicture}
